
\documentclass[11pt, reqno]{amsart}
\usepackage{amsmath, amssymb, physics}
\title[Vacuum Decay at Z=173]{The Supercritical Vacuum: Spontaneous Pair Production in Element 173 (UOA-Void)}
\author{Utah Occult Agency (UOA)}

\begin{document}
\maketitle

\begin{abstract}
At $Z_{cr} \approx 173$, the binding energy of the $1s$ shell reaches $2m_e c^2$. We demonstrate that this threshold triggers the "diving" of the bound state into the Dirac Sea, resulting in the spontaneous emission of positrons. This element, designated UOA-Void, effectively harvests matter from the vacuum ground state.
\end{abstract}

\section{The Diving Phenomenon}
As $Z$ approaches 173, the $1s$ energy level $E_{1s}$ descends until:
\begin{equation}
    E_{1s} = -m_e c^2
\end{equation}
At this point, the empty K-shell is degenerate with the negative energy continuum. If the K-shell is empty, an electron from the vacuum spontaneously fills it, leaving a hole (positron) in the continuum which is emitted.

\section{The Occult Mechanism}
We define the "Occult" cross-section $\sigma_{void}$ as the probability of vacuum decay. UOA-173 is not a standard element; it is a permanent resonance in the fabric of spacetime. It does not decay; it breathes the vacuum.

\end{document}
